% ============================================================
%  From PRGs to OWFs, hard-core predicates, from OWFs to PRGs
% ============================================================

\section{Hard-core predicates}

\Cref{thm:stretch} tells us that a PRG adding a single bit is enough. How do we build one? We will see that PRGs and one-way functions are \textbf{equivalent}: one exists if and only if the other does. Since nobody knows whether one-way functions exist, building a PRG is exactly as hard as building a OWF; in practice, real-world PRGs are \textit{heuristic} constructions, believed but not proved to be secure.

\subsection{From PRGs to one-way functions}

The easy direction first.

\begin{framedthm}[label={thm:prg-owf}]{From PRGs to OWFs}
  Let $G : \{0,1\}^n \to \{0,1\}^{2n}$ be a PRG. Then the function $f : \{0,1\}^{2n} \to \{0,1\}^{2n}$ defined by
  \[
    f(x_1 \concat x_2) = G(x_1) \qquad \text{with } x_1, x_2 \in \{0,1\}^n
  \]
  is a one-way function.
\end{framedthm}

\textbf{Intuition.} Whoever can invert $f$ can \textbf{recognize} the outputs of $G$: take the string, try to invert it, and recompute. On outputs of $G$ the inversion works with noticeable probability, while a uniform string is almost never an output of $G$ at all.

\begin{proof}
  $f$ is efficiently computable because $G$ is. Suppose, by contradiction, that a PPT adversary $\Adv$ inverts $f$ with a probability $\varepsilon(n)$ that is not negligible. Build the distinguisher $D$ that, on input $z \in \{0,1\}^{2n}$:
  \begin{enumerate}
    \item computes $x' \leftarrow \Adv(z)$;
    \item outputs $1$ if $f(x') = z$, and $0$ otherwise.
  \end{enumerate}
  \begin{itemize}
    \item \textbf{If $z \leftarrow G(\Uni{n})$}: since $f$ ignores the second half of its input, $z = G(x_1)$ with $x_1$ uniform is distributed exactly as $f(x)$ with $x$ uniform. So $D$ outputs $1$ exactly when $\Adv$ inverts, that is, with probability $\varepsilon(n)$.
    \item \textbf{If $z \leftarrow \Uni{2n}$}: $D$ outputs $1$ only if $z$ is in the image of $G$. The image contains at most $2^n$ strings out of $2^{2n}$, so this happens with probability at most $2^{-n}$.
  \end{itemize}
  The advantage of $D$ is at least $\varepsilon(n) - 2^{-n}$, which is not negligible (for large $n$ it is at least $\frac{\varepsilon(n)}{2}$). This contradicts the fact that $G$ is a PRG.
\end{proof}

\subsection{Hard-core predicates}

For the opposite direction we need a bit of the input that stays \textbf{secret} even when $f(x)$ is known.

\begin{frameddefn}[label={def:hardcore}]{Hard-core predicate}
  Let $f : \{0,1\}^n \to \{0,1\}^m$ be efficiently computable. A function $h : \{0,1\}^n \to \{0,1\}$ is a \textbf{hard-core predicate} for $f$ if for every PPT adversary $\Adv$:
  \[
    \Pr\big[\Adv(f(x)) = h(x) : x \sample \{0,1\}^n\big] \leq \frac{1}{2} + \negl(n)
  \]
\end{frameddefn}

In words: seeing $f(x)$, no efficient adversary can guess the bit $h(x)$ noticeably better than tossing a coin. An equivalent formulation, which is often more convenient, is
\[
  \big(f(x),\ h(x)\big) \cind \big(f(x),\ \Uni{1}\big) \qquad \text{with } x \sample \{0,1\}^n
\]
that is, $h(x)$ looks like a fresh random bit even next to $f(x)$.

Neither definition requires $f$ to be one-way, but intuitively a hard-core bit is what makes a function hard to invert. This raises three natural questions.

\paragraph{Question 1: is there a universal hard-core predicate, hard-core for every OWF?} \textbf{No.}

\begin{proof}
  Suppose that an efficiently computable $h$ were hard-core for every OWF. Take any OWF $f$ and build a ``bad'' function that \textbf{publishes} the bit $h(x)$ next to $f(x)$:
  \[
    f_{\mathrm{BAD}}(x) = f(x) \concat h(x)
  \]
  \begin{itemize}
    \item \textbf{$h$ is clearly not hard-core for $f_{\mathrm{BAD}}$}: the adversary that outputs the last bit of $f_{\mathrm{BAD}}(x)$ guesses $h(x)$ with probability $1$.
    \item \textbf{$f_{\mathrm{BAD}}$ is still one-way.} We prove it by reduction. Suppose that there is a PPT adversary $\Adv$ that inverts $f_{\mathrm{BAD}}$ with probability at least $\frac{1}{\poly(n)}$, that is, not negligible. We build a PPT adversary $B$ that inverts $f$, playing between the challenger $C$ of the one-wayness game of $f$ and the adversary $\Adv$:
  \end{itemize}

  \begin{center}
    \begin{tikzpicture}[x=1cm, y=1cm]
      \node (A) at (0,0) {$\Adv$};
      \node (B) at (4.5,0) {$B$};
      \node (C) at (9,0) {$C$};
      \node[lbl, align=left] at (9,-1.2) {$x \sample \{0,1\}^n$};
      \draw[msg] (9,-0.6) -- node[above] {$y = f(x)$} (4.5,-0.6);
      \node[lbl] at (4.5,-1.2) {$\tilde b \sample \{0,1\}$};
      \draw[msg] (4.5,-1.6) -- node[above] {$y \concat \tilde b$} (0,-1.6);
      \draw[msg] (0,-2.4) -- node[above] {$x'$} (4.5,-2.4);
      \draw[msg] (4.5,-2.4) -- node[above] {$x'$} (9,-2.4);
    \end{tikzpicture}
  \end{center}

  $B$ receives $y = f(x)$, \textbf{guesses} the missing bit with a coin $\tilde b$, gives $y \concat \tilde b$ to $\Adv$ and forwards its answer $x'$.
  \begin{itemize}
    \item With probability $\frac{1}{2}$ the guess is right, $\tilde b = h(x)$. The coin is independent of $x$, so in this case the input of $\Adv$ is distributed \textbf{exactly} as $f_{\mathrm{BAD}}(x)$ for a random $x$, and $\Adv$ finds $x'$ with $f_{\mathrm{BAD}}(x') = y \concat \tilde b$ with probability at least $\frac{1}{\poly(n)}$.
    \item $f_{\mathrm{BAD}}(x') = y \concat \tilde b$ implies in particular $f(x') = y$, so $B$ wins.
  \end{itemize}
  Therefore
  \[
    \Pr[B \text{ wins}] \geq \frac{1}{2} \cdot \Pr[\Adv \text{ wins}] \geq \frac{1}{2} \cdot \frac{1}{\poly(n)} = \frac{1}{\poly'(n)}
  \]
  which is not negligible, against the one-wayness of $f$. (The inequality is $\geq$ and not $=$ because $B$ may win even when the guess is wrong.)

  So $f_{\mathrm{BAD}}$ is a OWF for which $h$ is not hard-core, a contradiction.
\end{proof}

\textbf{The next best thing.} Since no single $h$ works for every $f$, we reverse the quantifiers: instead of ``$\exists\, h\ \forall f$'' we look for ``$\forall f\ \exists\, h$'', that is, a hard-core predicate tailored to each function. Question 3 shows that this is possible, up to a small modification of $f$.

\paragraph{Question 2: is every function with a hard-core predicate one-way?} \textbf{No.} Let $f$ delete the last bit, $f(x_1 \cdots x_n) = x_1 \cdots x_{n-1}$, and let $h(x) = x_n$ be the deleted bit. Since $x$ is uniform, $x_n$ is independent of $f(x)$, so $h$ is hard-core (no adversary can do better than $\frac{1}{2}$). But $f$ is trivially invertible: appending any bit to $y$ gives a valid preimage, so the adversary inverts with probability $1$.

\paragraph{Question 3: does every OWF have a hard-core predicate?} \textbf{Yes, up to a small modification.} This is the Goldreich--Levin theorem.

\begin{framedthm}[label={thm:gl}]{Goldreich--Levin}
  Let $f : \{0,1\}^n \to \{0,1\}^m$ be a OWF, and define $g : \{0,1\}^{2n} \to \{0,1\}^{m+n}$ as
  \[
    g(x, r) = f(x) \concat r
  \]
  Then $g$ is a OWF, and it has the hard-core predicate
  \[
    h(x, r) = \langle x, r \rangle \bmod 2 = x_1 r_1 \oplus x_2 r_2 \oplus \cdots \oplus x_n r_n
  \]
\end{framedthm}

How to read it:
\begin{itemize}
  \item \textbf{$r$ is public} $\to$ it appears in clear in the output $g(x, r)$, exactly like the seed of an extractor.
  \item \textbf{$h(x, r)$ is the XOR of a random subset of the bits of $x$} $\to$ the subset is selected by $r$: bit $x_i$ is included if and only if $r_i = 1$. It is the same family of functions as in the Alice and Bob extraction example.
  \item \textbf{$g$ is one-way} $\to$ the second half $r$ is in clear and carries no information about $x$, so inverting $g$ means inverting $f$.
\end{itemize}

\paragraph{Intuition: a perfect predictor would invert $f$.} Suppose there were a very good adversary $\Adv$ that, given $g(x, r) = f(x) \concat r$, always computes $h(x, r) = \langle x, r \rangle$, \textbf{for every} $r$. Then we could recover $x$ bit by bit, by choosing as $r$ the unit vectors $e_i = (0, \ldots, 0, 1, 0, \ldots, 0)$, with the $1$ in position $i$:
\[
  \langle x, e_1 \rangle = \langle x, 100\cdots0 \rangle = x_1, \qquad \ldots, \qquad \langle x, e_n \rangle = \langle x, 0\cdots01 \rangle = x_n
\]
So $n$ calls $\Adv(f(x) \concat e_i)$ would return all the bits of $x$, and $f$ would be inverted, against the hypothesis.

\paragraph{Why the real proof is harder.} The hypothesis only gives an adversary that guesses $h(x, r)$ with probability $\frac{1}{2} + \frac{1}{\poly(n)}$, and only \textbf{on average over a random $r$}. Such an adversary may be wrong on every unit vector $e_i$, which are very particular values of $r$. The full proof (omitted in class) has to work with random $r$ only.

\begin{framedobs}{Beyond the lecture: the self-correction trick}
  An intermediate case shows how to avoid the particular values $e_i$. Suppose that $\Adv$ guesses $\langle x, r \rangle$ with probability at least $\frac{3}{4} + \varepsilon$ over a random $r$. Since the inner product is linear,
  \[
    \langle x, r \rangle \oplus \langle x, r \oplus e_i \rangle = \langle x, e_i \rangle = x_i
  \]
  So, to compute $x_i$, choose a \textbf{random} $r$ and ask $\Adv$ both $\langle x, r \rangle$ and $\langle x, r \oplus e_i \rangle$. Both $r$ and $r \oplus e_i$ are uniformly distributed, so each answer is wrong with probability at most $\frac{1}{4} - \varepsilon$, and by the union bound both are right with probability at least $\frac{1}{2} + 2\varepsilon$. Repeating with many independent $r$ and taking the majority gives $x_i$ with high probability. The real proof, for probability $\frac{1}{2} + \varepsilon$, needs a cleverer way of choosing the queries.
\end{framedobs}

\subsection{From one-way permutations to PRGs}

Hard-core predicates are exactly what we need to build a PRG. The simplest case uses a \textbf{one-way permutation} (OWP): a OWF that is a bijection $f : \{0,1\}^n \to \{0,1\}^n$.

\begin{framedprop}[label={prop:owp-prg}]{From OWPs to PRGs}
  Let $f$ be a one-way permutation with hard-core predicate $h$. Then
  \[
    G(s) = f(s) \concat h(s)
  \]
  is a PRG with stretch $1$.
\end{framedprop}

\begin{proof}
  $G$ is efficiently computable because $f$ and $h$ are. For $s$ uniform:
  \begin{itemize}
    \item $f(s)$ is \textbf{uniform} over $\{0,1\}^n$, because a permutation of a uniform string is uniform;
    \item by the equivalent formulation of hard-core predicates, $h(s)$ looks like a fresh random bit even next to $f(s)$.
  \end{itemize}
  Therefore
  \[
    G(\Uni{n}) = \big(f(s),\, h(s)\big) \cind \big(f(s),\, \Uni{1}\big) = \big(\Uni{n},\, \Uni{1}\big) = \Uni{n+1} \qedhere
  \]
\end{proof}

Combined with Goldreich--Levin: if $f$ is a OWP, then $g(x, r) = f(x) \concat r$ is also a permutation, so $G(x, r) = f(x) \concat r \concat \langle x, r \rangle$ is a PRG with stretch $1$.

\subsection{From one-way functions to PRGs}

The proposition can be extended to \textbf{every} one-way function, not only permutations. The construction is much harder (it combines many evaluations of $f$ with their hard-core bits), and its proof was omitted.

\begin{framedthm}[label={thm:hill}]{Håstad, Impagliazzo, Levin, Luby}
  If a one-way function exists, then a PRG with stretch $1$ exists.
\end{framedthm}

Together with \Cref{thm:prg-owf} and \Cref{thm:stretch} this gives the complete picture:
\[
  \begin{aligned}
    \text{OWF} \;&\xrightarrow{\ \text{Goldreich--Levin, HILL}\ }\; \text{PRG with stretch } 1 \;\xrightarrow{\ \text{stretching}\ }\; \text{PRG with stretch } \poly(n) \\
    \text{PRG} \;&\xrightarrow{\ \text{\Cref{thm:prg-owf}}\ }\; \text{OWF}
  \end{aligned}
\]
\begin{framedcor}{}
  One-way functions exist if and only if pseudorandom generators exist.
\end{framedcor}
